\begin{figure}[hb!]
    \centering

    \hrule
    \vspace{0.9cm}

    % \pgfmathsetmacro{\mpolx}{\mstrong}
    % \pgfmathsetmacro{\mpoly}{1}
    % \pgfmathsetmacro{\qhat}{(\mpolx*\qpoly-\mpoly*\qpolx)/(\mpolx-\mpoly)}
    \pgfmathsetmacro{\xmin}{0}
    \pgfmathsetmacro{\xmax}{3}
    \pgfmathsetmacro{\delx}{0.6}
    
    \begin{tikzpicture}[scale=1.5]

        \begin{axis}[
            hide axis,
            xmin=\xmin, xmax=\xmax,
            ymin=\xmin, ymax=\xmax,
            width=3 cm,
            height=3 cm,
            at={(0,0)},
            scale only axis,
            clip=true,
            anchor=south west,
        ]
            \addplot[mygrey, no marks] 
            table [ x expr=\thisrowno{0}, 
                    y expr=\thisrowno{1}, 
                    col sep=comma,
                    unbounded coords=jump] 
            {figures/33/streamlines_greedy_expl.csv};

        \end{axis}

        \fill[CambridgeCoreBlue, opacity=0.1] (\xmax/2+\delx-0.5,\xmax/2+\delx-0.5) -- (\xmax,\xmax/2+\delx-0.5) -- (\xmax,\xmax) -- cycle ;
        \fill[CambridgeCoreBlue, opacity=0.1] (\xmax/2+\delx-0.5,\xmax/2+\delx-0.5) -- (\xmax/2+\delx-0.5, \xmax) -- (\xmax,\xmax) -- cycle ;
        \draw[CambridgeCoreBlue, thick] (\xmax/2+\delx-0.5,\xmax/2+\delx-0.5) rectangle (\xmax,\xmax);
        \draw[->] (\xmin, \xmin) -- (\xmax, \xmin) node[right, font=\small] {$q(s,\pol)$};
        \draw[->] (\xmin, \xmin) -- (\xmin, \xmax) node[left, font=\small] {$q(s,\pol_+)$};
        
        \draw[dotted] (\xmin, \xmin) -- (\xmax, \xmax);
        \draw[CambridgeCoreBlue, thick] (\xmax/2+\delx-0.5, \xmax/2+\delx-0.5) -- (\xmax, \xmax);

        \filldraw[fill=black, draw=black] (\xmax/2+\delx-0.5,\xmax/2+\delx/2-0.5) circle (1pt) node[below right, text=black] {$q_\pol$};
        \filldraw[fill=black, draw=black] (\xmax/2-\delx/2-0.5,\xmax/2-\delx-0.5) circle (1pt) node[below right, text=black] {$q_{\pol_+}$};

        \node at (2.53, 2.06) [color=CambridgeCoreBlue] {$q$};
        \node at (2.06, 2.53) [color=CambridgeCoreBlue] {$q_+$};

    \end{tikzpicture}
    \vspace{0.4cm}
    \caption{If the update distributions of mean-field MC-O-PI have a greedy bias and the solution transitions in state $s$ from a policy $\pol$ to another policy $\pol_+$ that satisfies the conditions in \eqref{eq: switch to worse}, the solution approaches $\qpol$ and $q_{\pol_+}$ from above. After the transition, the solution immediately moves back towards the boundary, introducing a sliding mode between the policies $\pol$ and $\pol_+$.}
    \label{fig: effect of greedy bias}
\end{figure}%
